% RESEARCH_DRAFT_SHA256: d06b1c9740133b19f7df8a45525d32fd765175a3cb8ae58a7d224db85d956310
% Layout and source-reference wrappers expanded for the web renderer.
\documentclass[11pt]{article}
\usepackage[margin=25mm]{geometry}
\usepackage{amsthm}
\usepackage{fontspec,unicode-math}
\setmainfont{Latin Modern Roman}
\setmathfont{Latin Modern Math}
\usepackage{amsmath,amssymb,mathtools,xurl,hyperref}
\hypersetup{colorlinks=true,urlcolor=blue,linkcolor=blue}
\setcounter{secnumdepth}{0}
\setlength{\parindent}{0pt}
\setlength{\parskip}{5pt}
\setlength{\emergencystretch}{3em}
\newcommand{\sref}[2]{\hyperlink{source:#1:#2}{[S#2]}}
\newcommand{\cataloguescope}[1]{\paragraph{Source record.}#1}
\newcommand{\floor}[1]{\left\lfloor#1\right\rfloor}
\newtheorem{proposition}{Proposition}
\newtheorem{lemma}[proposition]{Lemma}
\begin{document}
\paragraph{Research manuscript.} OpenAI \texttt{gpt-6-astra}, reasoning \texttt{ultra}; 11 October 2026.
\paragraph{Preserved source body.} The complete catalogue mathematical text and sources follow with layout and source links expanded, including source qualifications and its historical preparation notes. New research begins at the section headed Research attempt.

% UnsolvedMath ID 20002124; source catalogue code AIM-INFRASTRUCTURE-0024
\section{An access-validity workflow certificate for online assessment}\label{20002124}
{\small UnsolvedMath ID 20002124 · Computer Science}
\subsection{Definitions and mathematical statement}

\paragraph{Shared notation.}Standard number systems are $\mathbb N,\mathbb Z,\mathbb Q,\mathbb R,\mathbb C$. Specialized notation and index ranges follow the cited source and the definitions below.


Accessibility concerns whether learners with different sensory, motor, cognitive, or other access needs can perceive content, operate the interface, and communicate mathematical responses. An assistive technology is a tool supporting those activities. An accommodation adapts an activity to an identified learner need.
\paragraph{Record scope and qualifications.}
The catalogue title is retained as metadata. Named certificates, contracts, bounds or frameworks in that title are not asserted as results of the original source extract.
The source gives no tool, user requirements, accessibility standard, or verification procedure. These definitions provide context without prescribing an unmentioned legal or technical standard.

Catalogue screening limitation: Educational assessment, adoption, or study-design discussion; the supplied question does not state a mathematical target with defined assumptions and outcome.

The source snapshot labels this record partially solved; this is catalogue metadata rather than independent verification of a mathematical result or the current completion of a project.
\paragraph{Statement recorded in the supplied source.}
26. Accessibility of online tools for special needs \href{http://aimath.org/pastworkshops/oerassessmentproblems.pdf}{[S2]}
\paragraph{Source status.}
The supplied catalogue labels this source record partially solved; that metadata was checked against the supplied snapshot on 2026-10-11. The label is preserved for provenance and is not an independent verification of present open status or mathematical correctness.
\subsection{Short English statement}
Improve accessibility of online educational tools for learners with particular access needs.
\subsection{Sources}
\begin{itemize}
\item[S1] Ulam AI and the UnsolvedMath Contributors, UnsolvedMath: A Curated Collection of Open Mathematics Problems, record 20002124 (AIM-INFRASTRUCTURE-0024). Catalogue curation and metadata: CC BY 4.0. The source extract is retained; definitions, English summaries and qualifications are editorial additions.
\url{https://huggingface.co/datasets/ulamai/UnsolvedMath}.
\item[S2] AIM Workshop Problem Lists, source linked by catalogue record AIM-INFRASTRUCTURE-0024.
\url{http://aimath.org/pastworkshops/oerassessmentproblems.pdf}.
\end{itemize}
\paragraph{Preparation and review.}
The original source and UnsolvedMath contributors retain attribution for the source extract. Codex (GPT-6; exact serving revision not exposed) assisted with definitions, English summaries, scope notes and formatting. No new human mathematical authorship or independently verified solution is asserted. Permanent source ID 20002124 and code AIM-INFRASTRUCTURE-0024 checked against the complete supplied registry and current repository. Distinct upstream records, including repeated questions and broad or damaged extracts, are retained for complete coverage of the source CS category. A repeated record is not a new independent conjecture. Complete mathematical-equivalence and current-literature review has not been performed. Preparation and final review follow contribution prompts C and D, with source gaps disclosed. Structural and rendering checks do not establish a new theorem or open status.
\label{end:20002124}

\section{Research attempt}
\subsection{Literature, source context and scope}
The primary AIM notes describe a brainstorming session. The original excerpt and all catalogue qualifications remain above. They do not specify a mathematical theorem or a complete empirical study. The formal model below is an explicitly limited research route introduced for this attempt, not a reconstruction of a missing original conjecture. The supplied catalogue's earlier research summary was inspected as context; its named framework is not adopted as an established solution. The elementary result below is not claimed novel, and no educational effectiveness data were collected.

\subsection{Attempt: extending finite access tests across configurations}
The source names no accessibility standard or configuration set. I tested a conditional mathematical extension rule. Let configurations form a metric space $(X,d)$, and let a declared performance margin $f:X\to\mathbb R$ mean pass when $f\ge0$. Assume $f$ is $L$-Lipschitz and a finite tested set $S$ is an $\epsilon$-net: every configuration is within $\epsilon$ of a tested one.
\begin{proposition}
If every tested margin is at least $m$, then every configuration has margin at least $m-L\epsilon$. Without a known regularity bound, finitely many passing tests do not certify all untested configurations, even among continuous margin functions.
\end{proposition}
\begin{proof}
For any $x\in X$, choose $s\in S$ with $d(x,s)\le\epsilon$. The Lipschitz inequality gives $f(x)\ge f(s)-Ld(x,s)\ge m-L\epsilon$.

For the obstruction, suppose $x_*\notin S$. Since $S$ is finite, $a=\min_{s\in S}d(x_*,s)>0$. The constant function $f_0=1$ passes everywhere. The continuous function
\[
f_1(x)=1-2\max\{0,1-d(x,x_*)/a\}
\]
also equals one at every tested point, but equals $-1$ at $x_*$. Thus the finite observations cannot distinguish universal passing from failure. The second function has a finite Lipschitz constant $2/a$, showing why merely assuming continuity does not supply a quantitative certificate.
\end{proof}
\subsection{Verification and remaining gaps}
On $[0,1]$, testing $S=\{0,1/2,1\}$ gives net radius $1/4$. A known $L=2$ and tested margin $3/5$ imply a lower bound $1/10$. Without that bound, a spike centered at $1/4$ with $a=1/4$ passes all three tests and fails at its center. These are exact checks in an illustrative metric space. Actual access needs may be categorical and multidimensional; neither an appropriate metric nor a Lipschitz constant is provided by the source. This does not certify any real tool's accessibility or replace user testing and construct validity.

\section{References and source attribution}
American Institute of Mathematics, workshop discussion notes on open educational resources and assessment, \url{https://aimath.org/pastworkshops/oerassessmentproblems.pdf}. The 16-page primary PDF was opened on 11 October 2026 and its numbered brainstorming item was checked against the supplied record. The source supports the provenance and agenda, not the additional modeling assumptions or an empirical educational conclusion. Catalogue and contributor attribution are retained in the reproduced Sources section. The worked mathematics is proved here without reliance on an unverified external theorem.


\paragraph{Integration review.} Codex (GPT-6; exact serving revision not exposed) reviewed the written arguments and expanded layout and source-reference wrappers. Repository metadata and the production web renderer were checked using Python 3.13 and JavaScript with MathJax 3.2.2. This model-assisted review is not independent human proof verification; no PDF compilation or formal proof checking is asserted.

\section{Completion Estimate}
\noindent\textbf{COMPLETION: 0\%}
No substantiated new progress toward the entire original target is claimed. The worked results and failed route are delimited above; a known result, illustrative model, or auxiliary obstruction is not counted as original-target completion.

\section{Final status and attribution}
\textbf{Attempt status: unresolved.} This describes this research attempt, not a new ruling on the source catalogue's status or on external literature claims.
The mathematics and this self-audit were generated by OpenAI \texttt{gpt-6-astra}, reasoning \texttt{ultra}, on 11 October 2026. Codex (GPT-6; exact serving revision not exposed) orchestrated the batch and contributed formatting/repository review. No human mathematical authorship, independent human proof checking, or formal verification is claimed. The historical preparation paragraphs in the reproduced catalogue body concern that body only.

Contribution prompt A was applied to the complete definition and supplied source record, followed by the adversarial checks of prompt B. Local reproduction files are \texttt{batch008/checks.py}, \texttt{computations.json}, and the per-ID audit JSON. They were not public links at preparation time. Numerical checks are finite and do not prove unrestricted results. No LaTeX engine was found locally; compilation was not performed by this research agent.

\end{document}
